\chapter{AI-Based Multi-Label Food Spoilage Detection Using Image Analysis and IoT}
 Food spoilage is a significant challenge in the global food supply chain, leading to economic losses and health risks. Conventional spoilage detection methods are often time-consuming and subjective. Recent advances in artificial intelligence and Internet of Things (IoT) offer promising solutions for real-time and objective food quality monitoring. This project focuses on developing a multi-label classification system for food spoilage detection using image analysis and IoT integration. This chapter provides the background, literature review, objectives, methodology, and problem statement for the proposed study.

\section[Introduction]{\textbf{Introduction}}
Food spoilage is one of the major global problems affecting food quality, public health, and economic sustainability. Large quantities of fruits and vegetables are wasted every year due to microbial growth, environmental conditions, and improper storage practices. Traditional methods used for spoilage detection mainly rely on manual inspection based on visual appearance, smell, and texture changes. These approaches are subjective, inconsistent, and unsuitable for large-scale real-time monitoring applications. With the rapid advancement of Artificial Intelligence (AI), Deep Learning, Computer Vision, and Internet of Things (IoT) technologies, automated food quality monitoring systems have become an important area of research. AI-based image analysis techniques can identify spoilage indicators such as discoloration, mold formation, dark spots, and texture degradation more accurately and consistently than manual inspection methods.

This project proposes an AI-based multi-label food spoilage detection system using image analysis and IoT integration. Unlike conventional systems that classify food only as fresh or rotten, the proposed system categorizes fruits and vegetables into four spoilage stages: Fresh, Early Spoilage, Moderate Spoilage, and Severe Spoilage. Deep learning models such as CNN, MobileNetV2, ResNet50, and EfficientNetB0 are used for image classification, while IoT sensors such as MQ-135 and DHT11/DHT22 are used to monitor environmental factors including gas concentration, temperature, and humidity. The system also incorporates Explainable AI using Grad-CAM visualization to improve interpretability and transparency of predictions. A simple web-based interface is developed to provide real-time spoilage prediction and visualization, helping reduce food waste and improve food quality monitoring.

\section[Literature Review]{\textbf{Literature Review}}
Several researchers have explored the application of Artificial Intelligence, Deep Learning, Computer Vision, and Internet of Things (IoT) technologies for food spoilage detection and food quality monitoring. Conventional food inspection methods mainly depend on manual observation, which is often subjective, time-consuming, and inconsistent. To overcome these limitations, researchers have developed automated image-based systems capable of detecting spoilage indicators such as discoloration, mold formation, bruises, texture degradation, and surface damage in fruits and vegetables. Convolutional Neural Networks (CNNs) are among the most widely used approaches for food image classification because of their strong feature extraction and pattern recognition capabilities. Various studies have shown that CNN-based systems can achieve high classification accuracy for identifying spoiled food products under controlled conditions.

Recent advancements in transfer learning have further improved spoilage detection performance. Researchers have used pretrained architectures such as MobileNetV2, ResNet50, VGG16, EfficientNet, and DenseNet for feature extraction and classification tasks. These models are trained on large datasets such as ImageNet and fine-tuned for food spoilage applications, reducing training time and improving prediction accuracy even with limited datasets. Gupta et al. proposed an AI-based spoilage detection system using CNN models integrated with IoT sensors and ESP32 modules for real-time monitoring. Their work demonstrated that combining image analysis with environmental sensing improved detection reliability and enabled early spoilage prediction. Similarly, Mulla et al. developed a multimodal spoilage detection framework combining image processing and sensor data for improved food quality monitoring.

Researchers have also investigated advanced hybrid AI systems capable of performing multiple tasks simultaneously, including spoilage classification, freshness prediction, and shelf-life estimation. Transformer-based architectures combined with CNN and Long Short-Term Memory (LSTM) networks have shown promising performance in multi-task food analysis systems. Some studies used optical imaging and microscopy-based approaches combined with AI algorithms for bacterial detection and microbial spoilage analysis. These systems demonstrated high accuracy but often required expensive laboratory setups and controlled environments, limiting their practical implementation.

IoT-based monitoring systems have gained significant attention in recent years because environmental conditions play an important role in food deterioration. Temperature, humidity, and gas concentration influence microbial growth and spoilage progression. Several researchers integrated gas sensors such as MQ-135, ammonia sensors, and VOC sensors with microcontrollers such as Arduino, Raspberry Pi, and ESP32 for environmental monitoring. Sensor data collected from these devices were combined with machine learning algorithms to predict spoilage conditions and generate alerts. ESP32-based architectures are particularly popular because of their low cost, wireless connectivity, and ease of implementation in smart monitoring systems.

Explainable Artificial Intelligence (XAI) techniques have also become important in food spoilage detection research. Most deep learning systems operate as black-box models and do not provide information regarding how predictions are generated. To improve transparency and interpretability, researchers have used Grad-CAM and heatmap-based visualization methods to highlight image regions influencing classification results. These techniques help users understand whether the model focuses on meaningful spoilage indicators such as mold patches, dark spots, or texture variations during prediction.

Although existing research demonstrates significant progress in AI-based food spoilage detection, several limitations still remain. Most systems perform only binary classification such as fresh versus rotten and do not provide detailed information regarding spoilage progression. Many studies rely on small or non-diverse datasets, reducing the generalization capability of the models in real-world environments. Several systems also lack integration of environmental sensor data, real-time deployment capability, and Explainable AI features. Prototype-level implementations dominate current research, and practical user-friendly applications for household or small-scale food monitoring are still limited. Therefore, there is a need for a comprehensive system combining multi-label spoilage classification, IoT sensor fusion, Explainable AI, and real-time web-based deployment for efficient food quality monitoring and food waste reduction.

\section[Motivation]{\textbf{Motivation}}
Food spoilage is a major global issue that causes significant food waste and health risks. Fruits and vegetables are highly perishable and can spoil rapidly due to microbial growth, environmental conditions, and improper storage practices. Traditional spoilage detection methods mainly rely on manual inspection based on visual appearance, smell, and texture, which are often subjective, inconsistent, and time-consuming. These methods are not suitable for large-scale or real-time food quality monitoring applications.

Recent advancements in Artificial Intelligence (AI), Deep Learning, Computer Vision, and Internet of Things (IoT) technologies have enabled the development of automated food monitoring systems. AI-based image analysis techniques can identify spoilage indicators such as discoloration, mold formation, dark spots, and texture degradation more accurately than manual methods. IoT sensors can additionally monitor environmental factors such as temperature, humidity, and gas concentration, which strongly influence spoilage conditions.

Most existing food spoilage detection systems mainly perform binary classification such as fresh or rotten and provide limited information regarding spoilage progression. Many systems also lack real-time applicability, explainability, and environmental sensing integration. Therefore, this project is motivated by the need to develop a low-cost, intelligent, and user-friendly system capable of performing multi-label spoilage classification using AI and IoT technologies. The project also focuses on improving transparency and prediction reliability through Explainable AI and sensor fusion techniques for practical food quality monitoring and food waste reduction.

\section[Problem statement]{\textbf{Problem statement}}
Food spoilage leads to significant food waste, economic loss, and health risks due to the consumption of deteriorated food products. Traditional spoilage detection methods mainly rely on manual inspection based on visual appearance, odor, and texture changes, which are subjective, inconsistent, and unsuitable for real-time monitoring applications. Existing AI-based food spoilage detection systems primarily perform binary classification such as fresh or rotten and do not provide detailed information regarding the progression of spoilage stages. Many existing approaches also rely only on image analysis and fail to incorporate environmental factors such as temperature, humidity, and gas concentration, which play a major role in food deterioration.

Although deep learning models have shown promising results in image classification tasks, many current systems still lack Explainable AI integration, practical real-time implementation, and multimodal sensor fusion capability. In addition, limited datasets and lack of stage-wise spoilage classification reduce the effectiveness and real-world applicability of existing systems. Therefore, there is a need for an AI-based multi-label food spoilage detection system capable of classifying fruits and vegetables into Fresh, Early Spoilage, Moderate Spoilage, and Severe Spoilage stages using image analysis along with IoT-based environmental sensing. The proposed system aims to integrate deep learning, sensor fusion, Explainable AI, and a user-friendly interface to improve spoilage prediction accuracy, transparency, and practical food quality monitoring.

\section[Objectives]{\textbf{Objectives}}
The objectives of the project are:
\begin{enumerate}
\item To develop an AI-based system capable of classifying fruits and vegetables into multiple spoilage stages such as Fresh, Early Spoilage, Moderate Spoilage, and Severe Spoilage using deep learning techniques.
\item To integrate IoT-based environmental sensing using temperature, humidity, and gas sensors for improving spoilage detection reliability and real-time monitoring capability.
\item To develop a user-friendly web-based interface with Explainable AI visualization for real-time spoilage prediction, confidence analysis, and food quality monitoring.
\end{enumerate}

\section[Brief Methodology of the project]{\textbf{Brief Methodology of the project}}
The methodology adopted for the project involves dataset collection, preprocessing, deep learning model development, IoT integration, evaluation, and deployment. Initially, images of fruits and vegetables representing different spoilage stages are collected from publicly available datasets such as Kaggle and other open-source sources. The collected images are categorized into Fresh, Early Spoilage, Moderate Spoilage, and Severe Spoilage classes based on visible spoilage indicators such as discoloration, mold formation, dark spots, and texture degradation. The dataset is then preprocessed through resizing, normalization, cleaning, and augmentation techniques including rotation, flipping, zooming, and brightness adjustment to improve model performance and generalization capability.

Deep learning models such as CNN, MobileNetV2, ResNet50, and EfficientNetB0 are trained using transfer learning and fine-tuning approaches for spoilage classification. The dataset is divided into training, validation, and testing sets for performance evaluation. IoT sensors including MQ-135 gas sensors and DHT11/DHT22 temperature-humidity sensors are integrated using Arduino Uno or ESP32 microcontrollers to monitor environmental conditions influencing food spoilage. The image-based prediction confidence and sensor-based spoilage risk score are combined using a weighted fusion approach to generate the final spoilage prediction.

The trained models are evaluated using performance metrics such as accuracy, precision, recall, F1-score, confusion matrix, and ROC analysis. Explainable AI using Grad-CAM visualization is incorporated to highlight image regions influencing the prediction process. Finally, a simple Streamlit-based web interface is developed for real-time image upload, spoilage prediction, confidence visualization, and sensor monitoring.

\section[Assumptions made / Constraints of the project]{\textbf{Assumptions made / Constraints of the project}}
The food spoilage detection process is assumed to operate under controlled image acquisition and environmental monitoring conditions. Images collected from publicly available datasets are assumed to represent different spoilage stages accurately and consistently. The deep learning models are trained assuming sufficient dataset quality, balanced class distribution, and proper preprocessing. IoT sensor readings obtained from MQ-135 and DHT11/DHT22 sensors are assumed to provide reliable measurements of gas concentration, temperature, and humidity related to spoilage conditions.

The proposed system mainly focuses on fruits and vegetables selected for the dataset and may not generalize completely to all food categories. Environmental conditions such as lighting variations, camera quality, background noise, and image resolution may influence prediction accuracy. The system also assumes stable internet and hardware connectivity for real-time deployment and sensor integration. Due to limited project duration and computational resources, the dataset size, number of food categories, and model optimization scope are restricted. The project is also limited to selected spoilage stages and predefined environmental thresholds used for sensor fusion and prediction analysis.